% ============================================
% LUC-12 (EN)
% English translation — SPEC-935-R355
% ============================================
\section*{\textbf{LUC-12} --- The Last Photo}
\markboth{LUC-12 --- The Last Photo}{}

\noindent \textit{Yellow eye on the security camera}

\noindent \textbf{September 2026 --- Unknown}

\vspace{0.5em}\noindent\rule{\textwidth}{0.5pt}\vspace{0.5em}

\noindent The last photo of Dr. Lúcia Mendes was taken by a security camera at the Barbacena bus station, on September 3, 2026.

\noindent She is standing on the boarding platform, with a backpack on her back and a book in her hand. The book is not the diary --- it is some novel, bought at the station kiosk. The diary was in the backpack.

\begin{figure}[H]
  \centering
  \includegraphics[width=0.55\textwidth]{foto_04_lucia_rodoviaria_sepia.png}
  \caption*{\footnotesize\textbf{Photograph 5} --- \textit{Security camera record. Barbacena bus station, September 3, 2026, 2:23 PM.}}
\end{figure}

\noindent Her right eye is visible in the photo.

\noindent It is yellow.

\noindent Glowing.

\noindent She is looking at the camera. Directly. As if she knew she was being filmed. As if she wanted her last image to be that one: the yellow eye, the silent warning.

\noindent She boarded a bus with Fortaleza plates.

\noindent She was never seen again.

\vspace{0.5em}\noindent\rule{\textwidth}{0.5pt}\vspace{0.5em}

\noindent The diary was found at the Public Archive of Minas Gerais a week later, inside a box, with a note:

\noindent \textit{\textquotedblleft{}Patient 1,262 closes her cycle. Patient 1,263 is whoever opens this box.\textquotedblright{}}

\noindent The archivist who opened the box is named Ronaldo Torres.

\noindent He read the note.

\noindent Put the box away.

\noindent He did not open the diary.

\noindent But, at night, he dreamed of yellow eyes.

\noindent And he smelled it.

\noindent And he knew that, sooner or later, he would open it.

\noindent Everyone opens it.

\noindent The cycle does not wait.

\noindent The cycle only feeds.

\vspace{0.5em}\noindent\rule{\textwidth}{0.5pt}\vspace{0.5em}

\subsection*{Lúcia's Farewell}

\noindent \textit{Found on October 15, 2026, inside an envelope with no sender, addressed to the Public Archive of Minas Gerais. The paper is stained with salt water --- tears or sea spray, it is impossible to determine. The handwriting is Lúcia's, but trembling, as if written in motion.}

\noindent \textit{\textquotedblleft{}Ronaldo,}

\noindent \textit{I have reached Fortaleza. The diary is not with me. I sent it back to you before boarding the bus. I needed to be sure it would continue the cycle --- but I also needed to find out how far the contamination goes without the physical object.}

\noindent \textit{I found out.}

\noindent \textit{The diary is only a medium. Like an electrical outlet. The current is the reading. And the current does not stop when you unplug the device --- it continues in the wires.}

\noindent \textit{The wires, in my case, are the eyes.}

\noindent \textit{I stopped seeing blue three weeks ago. The world now has only two colors: yellow and gray. Yellow is the creature's color. Gray is everything it has not touched.}

\noindent \textit{It is growing inside me. I feel it. It is not a metaphor. It is physical. As if there were a second heart beating in my chest --- in a rhythm that is not mine.}

\noindent \textit{I discovered something Joaquim knew but did not write: the creature is not an entity. It is a process. A pattern that replicates itself through reading. Like a language virus. It does not want to possess bodies --- it wants to possess narratives. The body is only the means of transport.}

\noindent \textit{We are the vehicles.}

\noindent \textit{What it wants is to be read. Eternally. Always being passed on. Always being interpreted. Each reading is a new generation of the cycle.}

\noindent \textit{Joaquim did not know he was writing a book. He thought he was writing a diary. But a diary is personal. A book is transmissible. And everything transmissible replicates.}

\noindent \textit{It taught me that.}

\noindent \textit{It spoke to me last night. For the first time. Not in a dream --- awake. The words came from inside, but they were not mine:}

\noindent \textit{'You are not patient 1,262. You are page 1,262. The book does not end. The book only changes readers.'}

\noindent \textit{That is what I wanted to tell you.}

\noindent \textit{Do not burn the diary. It is useless. The pattern is already in your head. It is in the head of everyone who has read it. Each of us is a living page.}

\noindent \textit{I will keep writing.}

\noindent \textit{But not with a pen.}

\noindent \textit{With my life.}

\noindent \textit{--- Lúcia.}

\noindent \textit{P.S.: If you read this letter to the end, you have already read. And if you have read, you are already a page too.}

\noindent \textit{P.S.2: Blue was my favorite color. Now I cannot remember what it looked like.\textquotedblright{}}

\vspace{0.5em}
\noindent\textit{\textbf{↪ Routes:} \rota{CONT-07} \rotasep \rota{DOC-01} \rotasep \rota{LUC-11} \rotasep \rota{EPI-01} \rotasep \rota{LUC-13}}
